% !TeX program = xelatex
% ==========================================================================
% Пропоузал семестрового проєкту — курс «Штучний інтелект», УКУ.
% Це ЄДИНИЙ файл, який вам потрібно редагувати. Уся верстка й брендинг
% живуть у ucuproposal.cls — не змінюйте шрифти, поля, кеглі й кольори.
%
% Компіляція: XeLaTeX + Biber (Overleaf робить це автоматично).
% Локально:   latexmk (див. README.md).
% ==========================================================================
\documentclass[ukrainian]{ucuproposal}
% Опції класу (можна комбінувати через кому):
%   ukrainian | english  — мова документа й логотипу (за замовчуванням ukrainian)
%   slogan                — додає гасло УКУ в колонтитул титулки (за замовчуванням вимкнено)

% --------------------------------------------------------------------------
% Метадані титульної сторінки
% --------------------------------------------------------------------------
\projecttitle{Бенчмаркінг SLAM-алгоритмів на власному датасеті кампусу УКУ}
% \projecttagline{Одне речення про суть проєкту (опційно — можна видалити рядок)}

\teammember{Теодор Талан}
\teammember{Олександр Цибрівський}
\teammember{Тебенько Ярема}
% \teammember{Прізвище Ім'я} % третій учасник (опційно)
% \teammember[https://linkedin.com/in/...]{Прізвище Ім'я} % опційне клікабельне посилання на профіль (LinkedIn тощо)

\mentor{Косован Олександр} % якщо ментора немає — видаліть текст і залиште \mentor{}, рядок не друкується

\track{research} % research | engineering — не друкується на титулці, лише для внутрішньої категоризації

\repo{https://github.com/oToTak/SLAM-algorithms-proposal} % опційно; якщо репозиторію ще немає — залиште порожнім, рядок не друкується

\proposaldate{2026-09-20}

\begin{document}
\maketitlepage

\section{Проблема та мотивація}
% Що саме ви розв'язуєте і чому це важливо й цікаво в контексті ШІ?
% Орієнтовно 1 абзац.
SLAM - базова технологія автономної навігації UGV, але за зниженої
освітленості падає контраст, зростає шум і з'являється змаз рухом, через
що класичні детектори втрачають стабільні точки, а траєкторія дрейфує або
трекінг обривається. Наявні бенчмарки охоплюють переважно освітлені
приміщення \cite{burri2016euroc} чи денне водіння \cite{geiger2012kitti}.
Попередня робота клубу UCU Autonomous UGV на датасеті LaMARia
\cite{krishnan2025lamaria} показала, що RGB-only SLAM тримає трекінг у
темряві лише за наявності інерційної фузії та глобальної оптимізації, але
сама фіксує обмеження: частково синтетичне затемнення, ручна зйомка
замість рухомої платформи, лише міські сцени - і планує зібрати датасет
під UGV-домен \cite{basystyi2026lowlight}. Такий датасет зібрано. Ми
оцінюємо та порівнюємо сучасні SLAM-методи на \mbox{UCU Campus SLAM
Dataset v1}: вхід - 13 сиквенсів (кампус, Стрийський парк, mixed) із
low-light сценаріями, парами ,,прямо/зворот'' та GNSS ground truth; вихід
- оцінені траєкторії, метрики ATE/RPE під єдиним протоколом і аналіз
деградації під кроною дерев та в сутінках. У контексті ШІ робота
зіставляє класичні геометричні пайплайни з learning-based методами
(DPVO \cite{teed2023dpvo}, DPV-SLAM \cite{lipson2024dpvslam}) поза межами
їхнього тренувального розподілу й показує, чиї припущення порушуються
першими. Унікальна частина, виконувана в межах цього курсу, -
перенесення протоколу оцінювання на власний UGV-датасет і додавання
виміру релокалізації, відсутнього в попередній роботі.

\section{Огляд джерел та наявних рішень}
% Що вже зроблено у світі за цією темою? Спирайтесь на references.bib,
% цитуйте через \cite{}. Орієнтовно 3-5 джерел.
Візуальний SLAM поділяється на кілька парадигм із різними припущеннями про
сцену: feature-based ORB-SLAM3 \cite{campos2021orbslam3} потребує
повторюваних кутових точок і паралаксу; direct-метод DSO
\cite{engel2018dso} спирається на сталість яскравості; Kimera-VIO
\cite{rosinol2020kimera} реалізує оптимізаційну візуально-інерційну
одометрію, а OpenVINS \cite{geneva2020openvins} --- фільтровий підхід на
базі MSCKF; DPVO \cite{teed2023dpvo} і DPV-SLAM \cite{lipson2024dpvslam}
представляють learning-based напрям, другий --- з бекендом loop closure.
Оцінюють їх переважно на бенчмарках у сприятливих умовах: EuRoC
\cite{burri2016euroc} в приміщенні, KITTI \cite{geiger2012kitti} на денному
водінні, LaMARia \cite{krishnan2025lamaria} --- у міському масштабі з
ширшою варіацією освітлення та протоколом оцінювання на геореференсних
контрольних точках.
Найближча до нашої робота \cite{basystyi2026lowlight} оцінила всі шість
систем на LaMARia під єдиним протоколом ATE/RPE: Kimera-VIO єдина пройшла
всі сиквенси, поєднавши низький RPE зі зростаючим ATE через відсутність
loop closure; DPVO і DPV-SLAM не втрачали трекінг, але давали похибки
близько 100 м у темряві; ORB-SLAM3, DSO та OpenVINS провалювалися на
складніших сиквенсах. Ми успадковуємо її baseline, набір методів і
конвенцію звітування неповних прогонів; внесок --- перенесення протоколу на
реальні дані з рухомої платформи та перевірка релокалізації.

\section{Дані}
% Які набори даних, звідки, ліцензія. Якщо збираєте самі — як саме.
% Якщо дані не потрібні — поясніть чому.
У проєкті використовуватиметься \textbf{UCU Campus SLAM Dataset}, наданий
стейкхолдером UCU APPS \cite{ucu_slam_topic_2026}. Дані розміщені у
відкритому доступі на платформі Hugging Face, тому етап самостійного збору даних не передбачається.
\textbf{Обсяг і формат:} Набір містить 13 оцінювальних \cite{ucu_slam_topic_2026}
та 3 калібрувальні сиквенси. Усі дані доступні у форматах ROS 2 bag (MCAP) та TUM RGB-D. 
\textbf{Сенсори та середовище:} Дані записані за допомогою платформи
Clearpath Husky A200 і включають сенсори RGB-D, IMU та колісну
одометрію. Як еталонні вимірювання (Ground Truth) для оцінки
точності виступають GNSS-координати \cite{ucu_slam_topic_2026}. Записи
охоплюють структурований кампус УКУ та неструктурований Стрийський парк
в outdoor-умовах, включаючи специфічні сценарії зі слабким освітленням
(low-light) \cite{ucu_slam_topic_2026}.

\section{Підхід і методи}
% Які алгоритми/архітектури плануєте використати, чи спираєтесь на
% існуючі реалізації та як плануєте їх покращити.
У межах проєкту ми спиратимемося на існуючі open-source реалізації сучасних
SLAM-алгоритмів. Для репрезентативності ми відібрали системи, що
представляють ключові парадигми: feature-based, direct, filter-based
та learning-based \cite{ucu_slam_previous_team_2025}. 
Базовий стек для бенчмаркінгу включатиме: ORB-SLAM3, DSO, Kimera-VIO,
OpenVINS, DPVO та DPV-SLAM \cite{ucu_slam_previous_team_2025}. Наш підхід
полягає в інтеграції цих рішень та їх запуску на зібраних сиквенсах UCU Campus
Dataset \cite{ucu_slam_topic_2026}. Оскільки алгоритми тестуватимуться у складних
outdoor-умовах (зокрема low-light), частиною методології буде
адаптація та тюнінг конфігураційних параметрів кожного з алгоритмів
для максимізації їхньої стабільності перед проведенням фінальних
замірів \cite{ucu_slam_topic_2026, ucu_slam_previous_team_2025}.

\section{Оцінювання результатів}
% Метрики, baseline, спосіб порівняння, якісні очікування.
Для кількісної оцінки точності SLAM-алгоритмів ми використовуватимемо дві
стандартні метрики: абсолютну похибку траєкторії (ATE — Absolute
Trajectory Error) та відносну похибку позиціонування (RPE — Relative Pose Error) \cite{ucu_slam_topic_2026, ucu_slam_previous_team_2025}. 
Як еталонні дані (baseline / ground truth) для обчислення цих метрик
слугуватимуть вимірювання GNSS, які є частиною UCU Campus Dataset
\cite{ucu_slam_topic_2026}. Базовими показниками для порівняння
(baseline-алгоритмами) стануть результати класичних методів
(наприклад, ORB-SLAM3), які зазвичай демонструють нестабільність
та втрату трекінгу в умовах слабкого освітлення \cite{ucu_slam_previous_team_2025}.
\textbf{Якісні очікування:} 
Ми очікуємо, що алгоритми, які покладаються виключно на візуальні
дані, показуватимуть суттєву деградацію точності або відмови в
роботі у low-light сценаріях та неструктурованому середовищі
(наприклад, під кронами дерев у Стрийському парку)
\cite{ucu_slam_topic_2026, ucu_slam_previous_team_2025}. Натомість
системи з інерціальною ф'юзією (наприклад, Kimera-VIO) мають
продемонструвати вищу стабільність \cite{ucu_slam_previous_team_2025}.
Фінальним результатом стане готовий відтворюваний бенчмарк, придатний
для публікації на профільних воркшопах \cite{ucu_slam_topic_2026}.

\section{План і ризики}
% Етапи до midterm і до фінального захисту, ключові ризики та план Б.
\textbf{До 26.10.2026:}
\begin{itemize}
  \item розгортання всіх шести систем (ORB-SLAM3, DSO, Kimera-VIO, OpenVINS, DPVO, DPV-SLAM) - окремо;
  \item sanity-check на датасетах EuRoC/KITTI - звірка з опублікованими цифрами;
  \item єдиний eval-пайплайн ATE/RPE із синхронізацією GNSS ground truth;
  \item перший прогін на світлих сиквенсах - контрольна точка.
\end{itemize}

\textbf{До 21.12.2026:}
\begin{itemize}
  \item повний прогін на всіх 13 сиквенсах;
  \item фіксований протокол тюнінгу конфігурацій;
  \item запуск системи строгої релокалізації у відомих точках;
  \item аналіз деградації в гірших умовах;
  \item агрегація результатів.
\end{itemize}

\textbf{Ризики:}
\begin{enumerate}
  \item Повна втрата трекінгу в темряві і деградація самого GNSS ground
    truth у поганих умовах - враховувати такий прогін як неповний,
    ненадійні сегменти GNSS - окремо.
  \item Велике навантаження на залізо через комбінаторність методів,
    сиквенсів та ітерацій тюнінгу конфігурацій - тестити
    конфігурації на меншій підмножині сиквенсів.
  \item Релокалізацію підтримують не всі методи - для методів без
    підтримки порівнювати шлях туди-назад.
\end{enumerate}

% Приклад врізки з ключовими цифрами (опційно, можна видалити):
% \begin{keyfacts}
% \textbf{Ключові цифри:} 3 датасети \textperiodcentered\ 2 baseline-моделі \textperiodcentered\ 4 тижні до midterm
% \end{keyfacts}

\section{Contribution statement}
% Хто за що відповідає. Відсотки не обов'язкові — якщо хочете їх вказати,
% додайте в квадратних дужках: \contribution[50\%]{Прізвище Ім'я}{...}
\begin{contributions}
  \contribution[33\%]{Талан Теодор}{Спільна участь у повному циклі проєкту: підготовка даних, запуск SLAM-алгоритмів, обчислення метрик та написанн документації}
  \contribution[34\%]{Цибрівський Олександр}{Спільна участь у повному циклі проєкту: підготовка даних, запуск SLAM-алгоритмів, обчислення метрик та написанн документації}
  \contribution[33\%]{Тебенько Ярема}{Спільна участь у повному циклі проєкту: підготовка даних, запуск SLAM-алгоритмів, обчислення метрик та написанн документації}
\end{contributions}

% Примітка: декларація використання AI-інструментів (AI usage) до цього
% пропоузалу не входить — вона живе в README.md вашого репозиторію.

\ucuprintbibliography

\end{document}
